\chapter{Finance}

À partir des bilans de Tecfit pour les exercices 2022, 2023 et 2024, une analyse financière des ratios clés a été réalisée afin d’évaluer la stabilité financière de l’entreprise et sa capacité à mener à bien un projet de recherche d’envergure.

L’étude s’appuie sur dix ratios financiers issus de la liste de la Banque nationale de Belgique (BNB). Pour chacun d’entre eux, nous présentons brièvement son intérêt, son évolution au fil des années, ainsi qu’une comparaison avec un acteur du même secteur.

L’entreprise de référence retenue est Aqtor, spécialisée dans la fabrication de prothèses et d’orthèses. Active depuis 2018 et intégrée au groupe Eqwal Ability Group, Aqtor opère également en Belgique, dans la région de Gand. Ce choix se justifie par la proximité d’activité entre les deux sociétés et la pertinence de la comparaison sectorielle.

Les données 2023 d’Aqtor ont été utilisées, car elles offrent une base d’informations complète et récente, nécessaire pour assurer la fiabilité des comparaisons.

\section{Marge Nette sur vente}
Le premier ratio analysé est la marge nette sur ventes, également appelée taux de valeur ajoutée dans la nomenclature de la Banque nationale de Belgique (BNB). Ce ratio mesure la proportion du bénéfice d’exploitation par rapport au chiffre d’affaires, autrement dit, la part du chiffre d’affaires qui demeure après déduction de l’ensemble des charges opérationnelles. On a préféré utiliser la marge nette sur vente plutôt que la marge brute sur vente car les amortissements comptables ne tiennent pas compte de la totalité du matériel nécessaire à l’activité. Par exemple, Tecfit dispose d’une pompe à vide achetée d’occasion qui a une vingtaine d’année et qui peut tenir encore longtemps car peu utilisée. Son amortissement comptable reflète donc mal son utilisation réelle.  
L’évolution du ratio sur les trois dernières années est la suivante :

\begin{table}[h]
\centering
\begin{tabular}{c|c|c}
\textbf{2022} & \textbf{2023} & \textbf{2024} \\
\hline
0,29 & 0,33 & 0,28 \\
\end{tabular}
\end{table}

Ces chiffres indiquent qu’en 2024, pour 1euro de ventes, 24 centimes ont été transformés en bénéfice d’exploitation. Malgré une légère baisse par rapport à 2023, ce niveau reste globalement satisfaisant et nettement supérieur à celui du concurrent direct, dont la marge nette s’élève à 0,09.
Cette différence traduit une rentabilité opérationnelle solide et une bonne maîtrise des coûts, suggérant que l’entreprise parvient à exploiter efficacement ses ressources pour générer du profit. La hausse observée en 2023 s’explique par un bond du bénéfice d’exploitation de 57\%, plus important que celui du chiffre d’affaires (+40\%), liée à une optimisation des charges (diminution du compte 60/61 mais augmentation du CA entre 2022 et 2023).
En revanche, la baisse de 2024 mérite attention: elle résulte d’une augmentation du chiffre d’affaires (+15\%) mais une légère diminution du bénéfice d’exploitation (-1.2\%).

\subsection{Taux de valeur ajoutée}
Ce ratio mesure la proportion entre la marge brute d’exploitation et la somme de la marge brute d’exploitation, des achats de matières premières et des services externes. Il reflète la capacité de l’entreprise à créer de la valeur à partir des ressources consommées pour sa production. Un ratio plus élevé traduit une utilisation plus efficace des charges externes et des matières premières.
L’évolution du ratio pour Tecfit est la suivante :

\begin{table}[h]
\centering
\begin{tabular}{c|c|c}
\textbf{2022} & \textbf{2023} & \textbf{2024} \\
\hline
0,27 & 0,44 & 0,45 \\
\end{tabular}
\end{table}
On observe une progression significative entre 2022 et 2023, avec un taux de valeur ajoutée passant de 0,27 à 0,44. Cette hausse s’explique principalement par une augmentation du chiffre d’affaires (de € 792 505 à € 1 108 783) combinée à une baisse des charges liées aux achats et services externes (compte 60/61 : de € 733 451,64 à € 661 167). Par conséquent, la marge brut d’exploitation a nettement progressé (de € 265,061.00 à € 511,992.00).
En 2024, le ratio se maintient à un niveau stable (0,45), traduisant une bonne efficacité opérationnelle constante. On note cependant une hausse des charges d’exploitation (compte 60/61 : + € 68 644 par rapport à 2023). Le maintien du ratio vient d’une augmentation identique de la marge d’exploitation et du compte 60/61, ils ont tous les 2 augmenté d’approximativement de 15\%. 
Par comparaison, Aqtor affichait en 2023 un ratio de 0,47, approximativement identique à celui de Tecfit. Cela confirme que Tecfit gère efficacement ses matières premières et services externes, dégageant une marge proportionnellement identique à celle de plus gros concurrent que lui.

\subsection{Valeur ajoutée par personne occupée}
Ce ratio mesure le bénéfice d’exploitation généré par chaque employé. Il permet d’évaluer la productivité du personnel et d’identifier si les effectifs sont dimensionnés de manière optimale. Cet indicateur est particulièrement pertinent pour Tecfit, dont l’activité repose sur une production artisanale à forte intensité de main-d’œuvre, sans recours à l’automatisation. La création de valeur dépend donc directement de la performance et de l’efficacité de l’équipe.
Évolution du ratio
2022: 155 918,06€  |2023: 124 876,83€|2024: 133 174,59€
\begin{table}[h]
\centering
\begin{tabular}{c|c|c}
\textbf{2022} & \textbf{2023} & \textbf{2024} \\
\hline
 155 918,06€ & 124 876,83€ & 133 174,59€ \\
\end{tabular}
\end{table}
En 2022, le ratio atteint un niveau particulièrement élevé, principalement en raison d’un nombre réduit d’employés équivalents temps plein (1,7ETP). En 2023, ce nombre passe à 4,1ETP, soit une hausse de 140\%, alors que la marge d’exploitation a progressé de 93\%. Cette évolution explique la baisse du bénéfice par employé, sans pour autant traduire une baisse d’efficacité, on remarque au ratio 2 que les couts du compte 60/61 sont relativement constant alors que le chiffre d’affaires augmente, cette meilleure utilisation des ressources matériel se fait au prix d’un peu plus de travail pour la main d’œuvre. 
La légère augmentation observée en 2024 (+6.6\%) s’explique par une meilleure intégration des nombreux nouveaux employés, encore en phase de formation ou d’adaptation en 2023. 
Comparaison sectorielle
En comparaison, Aqtor présentait en 2023 une marge brute d’exploitation par employé de 173 562€, soit un peu moins de 40\% de plus que les employés de Tecfit. Cela met en évidence qu’il y a une possibilité d’amélioration pour Tecfit, bien qu’il ne puisse probablement pas égaler les performances d’Aqtor qui bénéficie d’une meilleure spécialisation de ses employés au vu de la taille de la structure. 
